\section{Introduction}
\label{sec:intro}

In October 2026, the Swiss helpline Die Dargebotene Hand (telephone 143) announced an AI chatbot for its online chat,
because only a fraction of chat requests could be answered by its volunteers \citep{srf2026dargebotene,keystone2026dargebotene}.
In crisis situations such as suicidal remarks, the bot follows a predefined procedure that includes an assessment of the
urgency of the situation and information about emergency services; in acute suicidality it stops and says that a human
is now needed \citep{keystone2026dargebotene,srf2026dargebotene}.
Such a bot therefore makes a judgement of urgency, and that judgement has to be communicated to people: to the
help-seeker, to the counsellor who takes over, and to the service that audits the conversation.
The sources do not describe how the bot makes this judgement.
In this paper we study one design in which urgency is expressed in \emph{words}, such as ``a crisis is
\word{likely}'', and ask whether those words mean the same thing to the machine and to the people it serves.

Probability words are vague.
Since Kent's work on estimative intelligence, it has been known that readers translate words such as \word{probable} or
\word{almost certain} into very different numbers \citep{kent1964words}, and survey studies have measured how wide and
how overlapping these translations are \citep{wallsten1986measuring,mosteller1990quantifying,budescu2014ipcc}.
Large language models (LLMs) now produce and read the same words at scale.
Recent studies find that LLMs reproduce the ordering of human readings but diverge for some words and contexts
\citep{belem2024perceptions,tang2026estimative,petridis2026unlikely}.
A bot that tells a counsellor that a crisis is \word{likely} transmits a number that the counsellor will reconstruct with
a human meaning.
If machine and human meanings differ, the same estimate can lead to a different decision.

\emph{Computing with Words} (CWW) offers a way to make such meanings explicit \citep{zadeh1996fuzzy}.
In Mendel's \emph{\Perc{}} (Per-C), words are encoded as interval type-2 fuzzy sets (IT2 FSs) whose footprint of
uncertainty records how much people disagree about a word, and these word models are collected in a \emph{codebook}
\citep{mendel2001perceptual,mendel2010perceptual}.
We use this architecture to compare a \emph{citizen codebook}, built from a large public survey of probability words, with an
LLM codebook, built by asking four local open-weight LLMs the same question many times.
We then insert both codebooks into the same crisis-line decision rule and count how often the choice of codebook changes
the action.

We ask three research questions, which follow a frozen analysis plan (Section~\ref{sec:design}):
\begin{description}
\item[RQ1] What do probability words mean to survey respondents, and how much do demographic groups disagree?
\item[RQ2] Do LLMs mean the same thing by these words as the respondents, and are LLM meanings narrower?
\item[RQ3] Does the mismatch between the citizen and the LLM codebook change the escalation decision of a crisis-line
\Perc{}?
\end{description}

The paper makes four contributions.
First, it proposes a \Perc{} for citizen-aligned risk words, with an encoder that builds IT2 word models from point
estimates via demographic group envelopes, an adaptation of the interval-based encoders of the CWW literature
(Section~\ref{sec:concept}).
Second, it reports a pre-registered comparison of the citizen and LLM codebooks for \WoneNTerms{} probability words,
including where LLM meanings collapse to a single number (Section~\ref{sec:results}).
Third, it quantifies how often the codebook choice flips the action of a crisis-line decision rule on real help-seeking
messages from public mental-health datasets (Section~\ref{sec:results}).
Fourth, it provides an open prototype, the word-meaning lab, that lets non-experts explore both codebooks and the
resulting decisions (Section~\ref{sec:prototype}).

Section~\ref{sec:background} introduces CWW, word encoding, verbal probability in people and LLMs, and the context of
cognitive cities and crisis lines.
Section~\ref{sec:concept} presents the concept, and Section~\ref{sec:design} the study design.
Section~\ref{sec:results} reports the results, Section~\ref{sec:prototype} the prototype,
Section~\ref{sec:discussion} discusses implications and limitations, and Section~\ref{sec:conclusion} concludes.
